\documentclass[12pt]{article}
% Research Paper Tracking Template

\usepackage[margin=1in]{geometry}
\usepackage{XCharter}
\usepackage{parskip}
\usepackage{enumitem}
\usepackage{booktabs}
\usepackage{tabularx}
\usepackage{array}
\usepackage{hyperref}
\usepackage{titlesec}
\usepackage{fancyhdr}

% ── Section formatting ──────────────────────────────────────────────────────
\titleformat{\section}{\fontsize{14}{16.8}\selectfont\bfseries}{}{0em}{}
\titlespacing*{\section}{0pt}{1em}{0.3em}

% ── Header / footer ─────────────────────────────────────────────────────────
\pagestyle{fancy}
\fancyhf{}
\rhead{\small Research --- \PaperYear}
\lhead{\small \CourseName}
\rfoot{\small \thepage}
\renewcommand{\headrulewidth}{0.4pt}
\setlength{\headsep}{12pt}

% ── Fill-in macros (edit for each paper) ────────────────────────────────────
\newcommand{\CourseName}{CS 6727 Cybersecurity Practicum}
\newcommand{\ProjectTitle}{Multi-Signature Authentication Web Proxy}
\newcommand{\AuthorName}{Ian Barish}
\newcommand{\PaperTitle}{[Paper Title]}
\newcommand{\PaperAuthors}{[Author Names]}
\newcommand{\PaperYear}{[Year]}
\newcommand{\PaperURL}{[URL or DOI]}

% ────────────────────────────────────────────────────────────────────────────
\begin{document}

% ── Title block ─────────────────────────────────────────────────────────────
\noindent Project: \ProjectTitle \hfill \AuthorName \\[0.2em]
\begin{center}
  {\large \textbf{\PaperTitle}} \\[0.3em]
  \PaperAuthors \ (\PaperYear)
\end{center}

\noindent\rule{\linewidth}{0.4pt}

% ── 1. Bibliographic Information ────────────────────────────────────────────
\section{Bibliographic Information}

\begin{itemize}[leftmargin=1.5em, itemsep=2pt]
  \item \textbf{Title:} \PaperTitle
  \item \textbf{Authors:} \PaperAuthors
  \item \textbf{Year:} \PaperYear
  \item \textbf{URL/DOI:} \PaperURL
\end{itemize}

% ── 2. Key Insight ──────────────────────────────────────────────────────────
\section{Key Insight}

[Summarize the paper's main contribution in 1--2 sentences. What is the core finding or innovation?]

% ── 3. Relevance to Project ─────────────────────────────────────────────────
\section{Relevance to Project}

[How does this paper relate to the multi-signature authentication web proxy project? Does it address authentication, web proxies, multi-person approval, security, usability, or another aspect of our work? Why is it worth knowing about?]

% ── 4. Challenges and Solutions ─────────────────────────────────────────────
\section{Challenges and Solutions}

\begin{itemize}[leftmargin=1.5em, itemsep=4pt]
  \item \textbf{Challenge 1:} [What problem did the authors identify or tackle?]

    \textit{Solution:} [How did they address it?]

  \item \textbf{Challenge 2:} [Another challenge they discuss...]

    \textit{Solution:} [How they addressed it...]

  \item \textbf{Limitation:} [Any limitations or open questions they mention?]
\end{itemize}

% ── 5. Impact on Project Approach ───────────────────────────────────────────
\section{Impact on Project Approach}

[How should this paper influence our project? Should we:
\begin{itemize}[leftmargin=1.5em, itemsep=2pt]
  \item Adopt or adapt a technique from this work?
  \item Avoid a pitfall they encountered?
  \item Reconsider our evaluation approach?
  \item Adjust our scope or deliverables?
  \item Update our timeline based on complexity insights?
\end{itemize}
Be specific about what changes, if any, we should make to our current plan.]

% ── 6. Notes ────────────────────────────────────────────────────────────────
\section{Additional Notes}

[Any other observations, quotes, related references, or thoughts worth preserving.]

\end{document}
